\section{Unified near-electrode model and numerical implementation}
\label{sec:methods}

\subsection{Resolving the sheath directly, rather than matching or
homogenizing it}
\label{sec:methods-strategy}

Section~\ref{sec:intro-electrode} distinguished three published strategies
for coupling a bulk-plasma closure to the electrode: solving one boundary
value problem with Poisson's equation over the whole gap
\cite{Almeida2008}; matching a non-equilibrium bulk solution directly to an
analytic sheath closure at an interface, valid strictly only for a
collision-dominated sheath \cite{Benilov2016}; and embedding the
unresolved sheath as an effective electrical conductivity in the
boundary-adjacent finite-volume cells of an otherwise quasi-neutral CFD
solution, $\sigma_{\rm eff} = \left[(1-I_g)/\sigma_{\rm child} +
I_g/\sigma\right]^{-1}$ with $\sigma_{\rm child}$ a Child-law space-charge-
limited conductivity \cite{LiangTrelles}. We use a fourth strategy,
implemented in the \textsc{PlasmeX} code: solve the same drift-diffusion--
Poisson system, equations~(\ref{eq:continuity})--(\ref{eq:poisson}), in the
bulk plasma and through the sheath alike, and use block-structured adaptive
mesh refinement (AMR), built on the AMReX framework, to refine the mesh at
the electrode down to a fraction of the local Debye length. No separate
sheath closure, matching condition, or homogenized conductivity is
introduced; the degree of local non-equilibrium is a property of the
converged solution, not a modelling choice made in advance. This is closest
in spirit to the unified modelling of \cite{Almeida2008}, generalized from
a specialized one-dimensional boundary-value solver to a general-purpose,
AMR finite-volume PDE code, and it should be contrasted directly with the
homogenized effective-conductivity shortcut of \cite{LiangTrelles}, which
is the more practical choice for genuinely three-dimensional geometries
where literal sub-Debye-length mesh refinement is not affordable.

\subsection{Transport closures}
\label{sec:methods-transport}

Electron mobility and diffusion coefficients are obtained from a
pre-tabulated, effectively Maxwellian two-term Boltzmann solution,
$\mu_eN(T_e)$ and $D_eN(T_e)$. This is verified, not assumed: across the
full tabulated range the ratio $(D_e/\mu_e)/T_e$ deviates from unity by
0.10\% rms (worst case 0.50\%), which is exactly the Einstein relation for
a Maxwellian electron energy distribution, and confirms that at the
0.9--10\% ionization fractions relevant here ($\nu_{ei}/\nu_{en}$ ranging
from 12 to 143 over $T_{\rm gas}=1000$--$12\,000$~K) electron--electron and
electron--ion collisions dominate every inelastic channel strongly enough
to Maxwellize the bulk electron population, exactly as required by
Section~\ref{sec:model-scaling}. The electron--ion Coulomb correction,
equation~(\ref{eq:coulomb}), reduces the neutral-collision-only
conductivity by one to three orders of magnitude and corrects its
qualitative trend (neutral-only conductivity rises with temperature only
because it divides by a falling neutral density; the corrected, Spitzer-like
conductivity, $\propto T_e^{3/2}/\ln\Lambda$, falls in the $\sim10^3$~S/m
range expected for LTE argon at 1~atm). A Saha-equilibrium cross-check on
the arc-core initial condition, $n_e^{\rm Saha} = 9.909\times10^{22}\,
\mathrm{m^{-3}}$ against the solved $n_e = 9.860\times10^{22}\,
\mathrm{m^{-3}}$, agrees to 0.5\%, confirming the bulk solution is a genuine
LTE state without that equilibrium being imposed. Ion transport uses the
closed-form two-temperature Wannier closure of
Section~\ref{sec:model-governing}; at the reduced field found at the
converged electrode cell of the sheath test case below (5976~Td), it gives
$\mu/\mu_0=0.16$, a factor-of-six reduction from the low-field mobility, and
recovers $\mu\to\mu_0$ in the low-field bulk.

\subsection{Electrode boundary conditions}
\label{sec:methods-bc}

At the electrode face, the electron flux is set to the standard kinetic
wall flux $\Gamma_e = \tfrac14 n_e \bar c_e$ evaluated on the (self-consistently
depleted) wall density, and the ion flux uses a Hagelaar-type drift-aware
closure, $\Gamma_i = \tfrac14 n_i\bar c_i + \max(\mu_i n_i E,0)$, that
reduces to the thermal flux at zero field and adds the outward drift
contribution when the local field drives ions into the wall. This is
structurally the same closure used for the counter-diffusing current in the
NLTE-sheath boundary condition of \cite{Benilov2016} (their $J_{CD}$,
equation~A.7) and for the back-diffusion current in \cite{LiangTrelles}
(their $J_{bd}$, equation~11): a Boltzmann-retarded thermal flux evaluated
at the wall density and local $T_e$. Two implementation choices in the
present model depart from a fully symmetric, field-aware treatment of both
species and are noted here as open items (Section~\ref{sec:results} /
\texttt{todo.txt}): the electron wall flux is currently evaluated with the
field-independent $\tfrac14 n_e\bar c_e$ form rather than a Hagelaar-type
field-aware form for electrons as well, and no secondary or thermionic
electron emission term is included at the electrode, which is normally the
dominant electron supply at a DC arc cathode and would need to be added
before any cathode-fall prediction from this model could be considered
complete.

\subsection{Numerical implementation}
\label{sec:methods-numerical}

The model is implemented in \textsc{PlasmeX}, a C++17 drift-diffusion
plasma code built on the AMReX block-structured AMR framework, currently
configured for a three-species (e$^-$, Ar$^+$, Ar) argon plasma. The
electrostatic potential is solved with AMReX's geometric multigrid solver
(MLMG) in general, and with Hypre's structured multigrid solver
(\texttt{HYPRE\_StructPFMG}) for the anisotropic two-dimensional
axisymmetric mesh used in production runs, since MLMG's composite
multigrid cycle was found not to converge on that mesh at any tested
aspect ratio, bottom solver, or coarsening depth. Species transport uses a
semi-implicit, operator-split scheme: drift fluxes are upwinded explicitly
and diffusive fluxes are backward-Euler in an implicit face-coefficient
form, $1 - \Delta t\,(e/\varepsilon_0)\mu_e n_e$, that removes the explicit
dielectric-relaxation time-step restriction of a fully explicit electron
drift treatment and gives roughly a $100\times$ larger stable time step at
equal accuracy, verified by matching both schemes at matched physical time
to $\sim10^{-6}$ relative agreement after a long integration from a common
checkpoint.

In one spatial dimension, both linear solves (the electrostatic Poisson
solve and the implicit species-diffusion solve) reduce to a tridiagonal
system on each AMR level, for which a direct solver is both faster and more
accurate than the general-purpose iterative multigrid solver: instrumented
at 93\% of a time step for a 512-cell problem, the iterative solver is
replaced by a direct Thomas-algorithm solve whose matrix entries are
recovered once, exactly, by probing the discrete operator, giving an
8.9$\times$ reduction in total time-step cost (65$\times$ and 13$\times$ on
the potential and species solves individually) with an exact rather than
iteratively converged linear-algebra answer. This extends, non-trivially,
to the full composite AMR hierarchy: refluxing at a coarse-fine interface
exactly cancels the covered coarse cell and replaces it with the fine-level
flux, and the coarse-to-fine ghost-cell interpolation used in one dimension
is an exact two-term closure over a single coarse cell,
$\mathrm{ghost} = \tfrac23 x_C + \tfrac13 x_{f_0}$ at a refinement ratio of
2 --- together these guarantee that the composite operator over the leaf
cells of a refined 1D hierarchy is itself exactly tridiagonal, so the same
direct solve applies unmodified across level boundaries. This direct solver
is what makes literal AMR refinement to sub-Debye-length cells in the
sheath computationally practical rather than prohibitively expensive; it is
restricted to one spatial dimension and a single MPI rank, so the
production two-dimensional axisymmetric configuration falls back to Hypre
as noted above.


